# P5 V2609 — Introduction Refinement (writing-shape session)
<!-- Raw pile: manuscript/V2609/Project5_GNN_Transformer_Antimalarial_main_V2609.tex §Introduction -->
<!-- Target: refined \section{Introduction} ready for drop-in replacement -->
<!-- Grounded so far: malaria burden, artemisinin resistance, PNA botanical names, co-evolution framing -->

The parasite is winning. Artemisinin-resistant \emph{Plasmodium falciparum} strains, first documented in Southeast Asia, have now been confirmed across sub-Saharan Africa --- the region that bears 94\% of the \num{608000} malaria deaths recorded in \num{2023} \citep{who2024malariareport,ariyey2014,ashley2018}. The pipeline of synthetic antimalarials has not kept pace. African medicinal plants have: \emph{Cryptolepis sanguinolenta}, \emph{Enantia chlorantha}, and \emph{Nauclea latifolia} produce indoloquinoline alkaloids, protoberberine scaffolds, and quassinoids shaped by millions of years of host--parasite co-evolution into structures no combinatorial library has replicated.

Screening this chemical space computationally reveals an immediate structural paradox. The \num{19836}-molecule panel examined here --- derived from \num{396} African natural products in ANPDB/AfroDb \citep{anpdb_2026} and expanded into a library of antimalarial candidates \citep{temgoua2026antimalarial} --- exhibits \qty{92.6}{\percent} whole-molecule novelty relative to standard synthetic screening libraries (Tanimoto $< 0.40$), yet retains \qty{69.3}{\percent} of privileged African ring cores, quantified by the African-Chemotype Structural Index (ACSI). The molecules are globally new; their scaffolds are not. This tension is not a curiosity: it is the structural signature that makes this panel a precise stress-test for molecular representation learning.

The structural density of these scaffolds --- elevated $F_{\mathrm{sp3}} \ge 0.45$, fused polycyclic cores, multiple stereocenters --- is precisely what breaks graph neural networks under scaffold separation. Standard message-passing architectures are bounded by the 1-Weisfeiler-Lehman (1-WL) expressivity test: they aggregate local atomic neighbourhoods but cannot resolve higher-order topological motifs such as bridged polycyclic cages or macrocyclic ring systems. On flat aromatic synthetics, this limitation is largely invisible. On African natural product scaffolds, it is the dominant failure mode.

The gap between fingerprint baselines and learned representations is not uniform across chemical space --- it is a function of structural distance from the training distribution. Nearest-neighbour Tanimoto decile stratification reveals that ECFP4--RF dominates only in the cold extrapolation tail ($D_{\mathrm{NN}} < 0.40$, deciles D1--D4), while deep models achieve near-parity in mid-range interpolation (D5--D8, $\Delta \in [0.001, 0.007]$). This distance-resolved picture, absent from prior benchmarks \citep{guo_ding_2026,benchmark_pretrained_2025}, is the analytical core of the Molecular Out-of-Distribution (MOOD) framework introduced here.

Fusing multi-scale persistent-homology images into a graph isomorphism network (GIN--TFP) injects global topological invariants --- zero-dimensional connectivity ($H_0$) and one-dimensional cycle persistence ($H_1$) --- that operate independently of local message passing. On polycyclic scaffolds (Rings $\ge 4$), this topological rescue recovers \qty{73}{\percent} of the cohort-level GNN deficit; at the aggregate scaffold split, it closes \qty{36}{\percent} of the gap. To exploit this distance-dependence prospectively, we formulate a Distance-Aware Conformal Triage Filter: a decision rule that routes each candidate to ECFP4--RF ($D_{\mathrm{NN}} < 0.40$) or GIN--TFP ($0.40 \le D_{\mathrm{NN}} \le 0.60$) based on its structural distance and expected calibration error, delivering a \qty{+14.2}{\percent} precision gain on the \num{22267}-compound ChEMBL transfer panel.

A fourth methodological contribution concerns molecular transformers specifically: enforcing per-fold pretrained weight restoration eliminates cross-fold information leakage that artificially inflates ChemBERTa performance estimates during cross-validation. Together, these four contributions --- distance-resolved mechanics, topological rescue, conformal triage, and leakage-free evaluation --- are operationalised through a centroid-based screening reduction that cuts compute overhead by \qty{99.3}{\percent} (\num{263000} to \num{1936} representative evaluations), making the full pipeline deployable on the HPC infrastructure available to malaria-endemic African research institutions without dependence on commercial cloud resources.

The central question this paper addresses is: \emph{under chemical-distribution shift, how do structural complexity and representation topology dictate the out-of-distribution extrapolation limits of molecular neural networks versus sparse bit-ensembles, and can distance-aware uncertainty triage dynamically optimise virtual screening performance?}

Random forests trained on ECFP fingerprints remain highly competitive on datasets of a few thousand molecules, and gains reported for GCN, GAT, or GIN architectures on large curated benchmarks often diminish significantly under scaffold splits \citep{guo_ding_2026,wu_molnet_2018,benchmark_pretrained_2025}. For natural products specifically, a \num{20}-fingerprint benchmark across $>\num{100000}$ compounds confirmed that no single encoding dominates \citep{boldini2024}. Topological data analysis introduces multi-scale geometric invariants through persistent homology \citep{tda_review_2025}, and tensor-network embeddings compress scaffold representations via Tucker decomposition \citep{tn_efficient_2026}; neither family has been systematically evaluated inside a graph fusion architecture under rigorous scaffold-held-out testing on an antimalarial natural product panel --- the gap this study closes.
